\section[Cote 72 : réalisations par pseudo-réflexions d'un polytope régulier]{Cote 72 : réalisations par pseudo-réflexions d'un polytope régulier\protect\verifie{Folder72}}\label{sec:72}

La cote 72, \emph{Polyèdres réguliers et hyperpolyèdres réguliers (calculs en dimension quelconque)}, que l'inventaire date \q{[à partir de 1978]}, étudie les réalisations du polytope régulier universel de rang $n+1$ par des pseudo-réflexions $\tau_i = \mathrm{id} + a'_i \otimes a_i$ ($0 \leq i \leq n$) d'un module $E$ de rang $n+2$ sur une base quelconque, les $\tau_i$ non consécutifs commutant. Avec le point de départ $s_0$, $s_{i+1} = \tau_i s_i$ et $a_i = s_{i+1} - s_i$, les accouplements sont $\langle a_{i-1}, a'_i\rangle = 1$, $\langle a_i, a'_i\rangle = \mu_i = \lambda_i - 1$, $\langle a_{i+1}, a'_i\rangle = \beta_i$, les autres nuls. On travaille ici dans la base $(s_0, a_0, \dots, a_n)$ de $E = R^{n+2}$, $R$ un anneau commutatif ; $\pi$ est la première coordonnée et $\tilde c$ l'unique vecteur annulé par tous les $a'_i$ dont la coordonnée sur $a_n$ vaut~$1$, obtenu par substitution remontante.

Page~37, $\mathcal U_\rho = \operatorname{Spec} \mathcal O_S[T]_{(1+\rho T)}$, de loi $\lambda * \lambda' = \lambda + \lambda' + \rho\lambda\lambda'$, opère par $u_\lambda = \mathrm{id} + \lambda\, \pi \otimes \tilde c$, où $\rho = \pi(\tilde c)$ ; le \q{Théorème} de la page~38 affirme sans démonstration que $\mathbb G_m \times \mathcal U_\rho \to \operatorname{Aut}(E, (\tau_i))$ est un isomorphisme et que $D_0 \smallsetminus \{c, d\}$, $D_0 = H_1 \cap \dots \cap H_n$, est un torseur sous $\mathcal U_\rho$ trivialisé par $s_0$ (\lot{72}{2}{38}). Les pages~87 et~88 le reprennent sous forme projective.

\begin{theoreme}[D'après les pages 37 et 38, en coordonnées, sur les points\piste{72-automorphisms-deform-gm-to-ga}]
Les $\lambda \in R$ avec $1 + \rho\lambda$ inversible forment un groupe commutatif $\mathcal U_\rho(R)$ ; $\lambda \mapsto 1 + \rho\lambda$ est surjectif sur les inversibles de toute $R$-algèbre si et seulement si $\rho$ est inversible, et $\mathcal U_0 = \mathbb G_a$. On a $u_\lambda u_{\lambda'} = u_{\lambda * \lambda'}$, $\det u_\lambda = 1 + \rho\lambda$, et un endomorphisme de $E$ commute à tous les $\tau_i$ si et seulement s'il s'écrit $c + t\, \pi \otimes \tilde c$, avec $(c, t)$ unique ; il est inversible si et seulement si $c$ et $c + t\rho$ le sont. D'où $R^\times \times \mathcal U_\rho(R) \simeq \operatorname{Aut}(E, (\tau_i))$, $(c, \lambda) \mapsto c\, u_\lambda$. Les points de $D_0$ sont les $a s_0 + t \tilde c$, et ceux qui évitent $c$ et $d$, normalisés par $\langle s, a'_0\rangle = 1$, sont exactement les $u_\lambda s_0 = s_0 + \lambda \tilde c$, chacun pour un seul $\lambda$, avec $u_\mu(u_\lambda s_0) = u_{\lambda * \mu} s_0$.
\end{theoreme}

\noindent L'énoncé tient sur les points, pour tout anneau : voir \texttt{forall\_surjective\_iff}, \texttt{u\_mul}, \texttt{det\_u}, \texttt{commute\_iff}, \texttt{isUnit\_iff}, \texttt{centraliserEquiv}, \texttt{mem\_startPoints\_iff} et \texttt{u\_e0\_injective} dans \href{https://github.com/Commutator-IO/grothendieck-project/blob/main/lean/Grothendieck/Folder72.lean}{\texttt{Folder72.lean}}. Le test de la réciproque est l'anneau des nombres duaux, où $1 + \varepsilon$ n'est atteint que si $\rho$ est inversible. Pour $n = 1$, $\rho = \mu_0\mu_1 - \beta_0$, soit $4 - \beta_0$ pour des réflexions, comme à la page~41.

Aux pages~79 à~82, sur un corps $k$, la cote décrit les sous-espaces invariants (\lot{72}{4}{79}). Avec $Z'_i = H_0 \cap \dots \cap H_i$ et $Z^*_i = \operatorname{Lin}(h_{i+1}, \dots, h_n)$, $h_i = k a_i$ et $H_i = \operatorname{Ker} a'_i$ :

\begin{theoreme}[Grothendieck, page 79\piste{72-invariant-subspaces-string}]
Sur un corps, les $\mu_i$ quelconques (caractéristique~$2$ comprise), un sous-espace de $E$ est stable par tous les $\tau_i$ si et seulement s'il est $0$, $E$, la droite $k\tilde c$ du centre, l'hyperplan $C$ d'équation $\pi = 0$, ou bien $Z'_i$ ou $Z^*_i$ pour un $i < n$ tel que $\beta_i = 0$.
\end{theoreme}

\noindent L'énoncé tient tel quel : voir \texttt{invariant\_iff} dans \href{https://github.com/Commutator-IO/grothendieck-project/blob/main/lean/Grothendieck/Folder72.lean}{\texttt{Folder72.lean}}, avec le corollaire~(c) de la page~80 (\texttt{invariant\_iff\_of\_ne\_zero}). Pour $i < n$, $Z^*_i \subset Z'_i$ si et seulement si $\beta_i = 0$ (\texttt{Zs\_le\_Zp\_iff}) : c'est le \q{$\in$} de la marge de la page~79 que la lecture retient, non le \q{$\notin$} de la page~81.

Ne sont pas démontrés ici : le schéma en groupes $\mathcal U_\rho$ et sa lissité, l'énoncé projectif sur une base (automorphismes de $X = \mathbf P(E)$ commutant aux $\tau_i$ projectifs), le torseur des drapeaux de la page~87, les comptes des pages~79 et~80 et l'irréductibilité sur $C$ des pages~80 à~82.
